% Gesamttest «Wellen» — Quelle des PDFs (python3 scripts/build-lp-pdf.py wellen).
% Musterlösung und Raster: bewertungspaket.tex. Alle Werte mit python3 nachgerechnet (09.10.2026):
% G1 T = 2.4 s, A = 10 cm, λ = 3.75 m/s · 2.4 s = 9.0 m (Tiefwasser: 1.56 · 2.4² ≈ 9.0 m; Wellenhöhe 0.2 m ≪ λ/7),
%    60 s / 2.4 s = 25 Berge pro Minute; G2 t = 102 m / 340 m/s = 0.30 s, λ = 340/170 = 2.0 m, λ_W = 1500/170 ≈ 8.8 m
%    (Gegenweg 1500/8.8 ≈ 170 Hz); G3 c = 1.4 m · 2.0 Hz = 2.8 m/s, t = 8.4/2.8 = 3.0 s; G4 λ = 340/50 000 = 6.8 mm,
%    λ = 3.00e8/2.4e10 = 1.25 cm; G5 f = 3.00e8/455e-9 ≈ 6.59e14 Hz, 637/455 = 1.4, 637 nm/0.4 ≈ 1590 nm;
%    G6 λ = 3.00e8/2.5e13 = 12 µm.
% Fassung 1.5 nach der Schlussprüfung: G2 Schall geht aus der Luft ins Wasser, Wellenlänge mit zwei Unterwassermikrofonen
%    gemessen; G1 «Sturmwellen» auf dem See.
% Fassung 1.4 nach der vierten Prüfung: G2 prüft die Wellenlänge im anderen Medium (statt s/λ), G4 mit 50 kHz.
% Fassung 1.3 nach der dritten Prüfung: jede Teilaufgabe gegen Kapitelaufgaben, Kontrollfragen, Clipprobleme,
% Leistenaufgaben und Übungen abgeglichen (Tabelle in scripts/lp/wellen/README.md, «Prüfung Runde 3»).
\documentclass[11pt]{article}
\usepackage{../lp-druck}
\renewcommand{\lpname}{Wellen}
\renewcommand{\lpteilgebiet}{6.1}
\renewcommand{\lpfuss}{Gesamttest · 25 Punkte}
\begin{document}
\lpkopf{Gesamttest}{%
  \vspace{4pt}\noindent\begin{tabularx}{\linewidth}{@{}X X@{}}
  Code (kein Name): \hrulefill & Punkte: \hrulefill\ / 25
  \end{tabularx}}

\begin{lpinfo}{Anleitung}
\textbf{Zeit:} rund 40 Minuten. \textbf{Hilfsmittel:} Taschenrechner und Formelsammlung, in allen Teilen.
Schreib zu jeder Aufgabe zuerst die Formel, dann die Werte \textbf{mit Einheiten}. Punkte gibt es für den Weg,
für das Ergebnis mit Einheit, fürs Ablesen und Zeichnen und für Begründungen in eigenen Worten. Rechne mit
\mbox{$c = 3.00 \cdot 10^{8}\;\text{m/s}$} (Licht und alle elektromagnetischen Wellen) und
\mbox{$c_\text{Luft} = 340\;\text{m/s}$} (Schall).
Danach mit dem \emph{Bewertungspaket} bewerten (Scan ohne Namen und Standort).
\end{lpinfo}

\lpteil{Teil A · Wellen beschreiben und berechnen}{Kapitel 1, 2 und 3 · 14 Punkte}

\begin{lpaufgabe}{G1}{Sturmwellen am Badefloss}{5 P}
Draussen auf dem See schaukelt ein Badefloss auf den Wellen eines Sturms; das Diagramm links zeigt seine Höhe über der Zeit. Die Wellen laufen mit \(3.75\;\text{m/s}\).\par\smallskip
\begin{center}\begin{tikzpicture}[x=0.95cm,y=0.085cm,>=latex]
\foreach \x in {0,0.2,...,6.01} \draw[lplinie!70,very thin] (\x,-15) -- (\x,15);
\foreach \y in {-15,-10,...,15} \draw[lplinie!70,very thin] (0,\y) -- (6,\y);
\foreach \x in {1,2,...,6} {\draw[lplinie] (\x,-15) -- (\x,15); \node[below] at (\x,-15) {\small \x};}
\foreach \y in {-10,10} {\draw[lplinie] (0,\y) -- (6,\y); \node[left] at (0,\y) {\small $\y$};}
\node[left] at (0,0) {\small 0};
\draw[->] (0,0) -- (6.15,0) node[right] {\small \(t\) in s};
\draw[->] (0,-15) -- (0,18) node[right] {\small \(y\) in cm};
\draw[very thick,lpbernstein,domain=0:6,samples=241] plot (\x,{10*sin(150*\x)});
\end{tikzpicture}\hspace{3mm}\begin{tikzpicture}[x=0.28cm,y=0.085cm,>=latex]
\foreach \x in {0,1,...,20} \draw[lplinie!70,very thin] (\x,-15) -- (\x,15);
\foreach \y in {-15,-10,...,15} \draw[lplinie!70,very thin] (0,\y) -- (20,\y);
\foreach \x in {5,10,15,20} {\draw[lplinie] (\x,-15) -- (\x,15); \node[below] at (\x,-15) {\small \x};}
\foreach \y in {-10,10} {\draw[lplinie] (0,\y) -- (20,\y); \node[left] at (0,\y) {\small $\y$};}
\node[left] at (0,0) {\small 0};
\draw[->] (0,0) -- (20.5,0) node[right] {\small \(s\) in m};
\draw[->] (0,-15) -- (0,18) node[right] {\small \(y\) in cm};
\end{tikzpicture}\end{center}
\begin{enumerate}[label=\alph*),leftmargin=*,itemsep=2pt,topsep=2pt]
\item Lies im Diagramm links die Amplitude und die Periode ab.
\item Berechne die Wellenlänge.
\item Zeichne im Raster rechts das Momentbild der Wasseroberfläche von \(s = 0\) bis \(s = 20\;\text{m}\), zu einem Zeitpunkt, an dem bei \(s = 0\) ein Wellenberg ist.
\item Wie viele Wellenberge treffen in einer Minute auf das Floss?
\end{enumerate}
\lpplatz{8}
\end{lpaufgabe}

\begin{lpaufgabe}{G2}{Konzert am See}{5 P}
Mia steht \(102\;\text{m}\) von der Bühne eines Konzerts am Seeufer entfernt. Gerade spielt die Band einen Ton von \(170\;\text{Hz}\). Ein kleiner Teil des Schalls dringt ins Wasser des Sees ein (Schall in Wasser: \(1500\;\text{m/s}\)).
\begin{enumerate}[label=\alph*),leftmargin=*,itemsep=2pt,topsep=2pt]
\item Wie lange braucht der Ton von der Bühne zu Mia?
\item Wie lang ist die Welle dieses Tons in der Luft?
\item Zwei Unterwassermikrofone im See werden in Laufrichtung des Schalls auseinandergeschoben. Für diesen Ton schwingen ihre Signale bei einem Abstand von \(8.8\;\text{m}\) erstmals wieder im gleichen Takt. Prüfe mit einer Rechnung, ob \(8.8\;\text{m}\) die Wellenlänge im Wasser sein kann.
\item Ein Taucher sagt: «Weil die Welle im Wasser länger ist, höre ich hier unten einen tieferen Ton.» Beurteile die Aussage.
\end{enumerate}
\lpplatz{8}
\end{lpaufgabe}

\begin{lpaufgabe}{G3}{Die Schraubenfeder}{4 P}
Eine lange Schraubenfeder liegt gestreckt auf dem Boden. Du bewegst ihr Ende hin und her.
\begin{enumerate}[label=\alph*),leftmargin=*,itemsep=2pt,topsep=2pt]
\item Einmal bewegst du das Ende seitlich, einmal in Richtung der Feder. Welche Wellenart entsteht jeweils? Begründe.
\item Jan sagt: «Deine Hand bewegt alle Windungen direkt mit; darum bewegen sich alle Windungen im selben Augenblick.» Beurteile die Aussage und erkläre, wie die Bewegung tatsächlich weitergegeben wird.
\item In der Mitte der Feder klebt ein Punkt. Wo ist er, nachdem die Störung vorbeigelaufen ist? Begründe damit, was die Störung transportiert.
\item Du bewegst das Ende mit \(2.0\;\text{Hz}\); die Wellen auf der Feder sind \(1.4\;\text{m}\) lang. Wie lange braucht ein Wellenberg bis ans Ende der \(8.4\;\text{m}\) langen Feder?
\end{enumerate}
\lpplatz{8}
\end{lpaufgabe}

\lpteil{Teil B · Elektromagnetische Wellen}{Kapitel 4 und 5 · 8 Punkte}

\begin{lpaufgabe}{G4}{Saugroboter und Schiebetür}{4 P}
Ein Saugroboter misst Abstände mit Ultraschall von \(50\;\text{kHz}\). Der Radarsensor einer Schiebetür sendet elektromagnetische Wellen mit \(24\;\text{GHz}\).
\begin{enumerate}[label=\alph*),leftmargin=*,itemsep=2pt,topsep=2pt]
\item Berechne beide Wellenlängen in Luft. In welchem Bereich des Spektrums liegen die Wellen des Radarsensors?
\item Die Frequenz des Radars ist 480\,000-mal so gross wie die des Ultraschalls. Erkläre, warum seine Welle trotzdem länger ist.
\item Erkläre, was bei der Radarwelle schwingt, und warum der Radarsensor darum auch durch ein Vakuum hindurch arbeiten könnte, der Ultraschall des Roboters aber nicht.
\end{enumerate}
\lpplatz{9}
\end{lpaufgabe}

\begin{lpaufgabe}{G5}{Zwei Farben aus einem Gas}{4 P}
Ein Gas aus gleichen Atomen leuchtet nur in zwei Farben: blau mit \(455\;\text{nm}\) und rot mit \(637\;\text{nm}\). Beim roten Licht springt ein Elektron von \(E_2\) auf \(E_1\), beim blauen von \(E_3\) auf \(E_1\).
\begin{enumerate}[label=\alph*),leftmargin=*,itemsep=2pt,topsep=2pt]
\item Welche Frequenz hat das blaue Licht?
\item Wie viel Mal grösser ist die Stufe \(E_3 - E_1\) als die Stufe \(E_2 - E_1\)? Rechne mit \(E = h \cdot f\).
\item Ein Elektron sitzt auf \(E_2\). Welche Wellenlänge muss ein Photon haben, das es genau auf \(E_3\) hebt? In welchem Bereich liegt es?
\item Aus solchen Atomen wird ein roter Laser gebaut. Erkläre, wie dort aus einem Photon viele gleiche werden, und wozu der Laser das Pumpen und die zwei Spiegel braucht.
\end{enumerate}
\lpplatz{9}
\end{lpaufgabe}

\lpteil{Teil C · Treibhauseffekt}{Kapitel 6 · 3 Punkte}

\begin{lpaufgabe}{G6}{Gase in der Atmosphäre}{3 P}
\begin{enumerate}[label=\alph*),leftmargin=*,itemsep=2pt,topsep=2pt]
\item Gib für Stickstoff, Wasserdampf und Ozon an, welche Strahlung jedes Gas aufnimmt: keine nennenswerte, das Ultraviolett der Sonne, die Wärmestrahlung des Bodens — auch mehreres ist möglich.
\item Begründe, warum mehr Kohlendioxid in der Luft den Boden im Mittel wärmer macht.
\item Welche Wellenlänge hat Infrarot mit \(2.5 \cdot 10^{13}\;\text{Hz}\)? Liegt sie im «Fenster» der Atmosphäre?
\end{enumerate}
\lpplatz{10}
\end{lpaufgabe}

\vfill
{\small\color{lpgrau}Bewerten: mit dem Bewertungspaket (Musterlösung, Punkteraster, Auftrag an eine KI) — erst nach dem Lösen öffnen.}
\end{document}
